\section{Benefícios esperados e desafios}

% \subsection{Contribuições esperadas}

% Espera-se que o presente projeto contribua para o avanço do estudo do PCmin, ainda pouco explorado na literatura, especialmente sob a perspectiva de formulações exatas e métodos de solução para instâncias de maior porte. Entre as principais contribuições esperadas, destaca-se o aperfeiçoamento das formulações matemáticas existentes, por meio da investigação de novas restrições válidas, eliminação de simetrias, redução de redundâncias e possíveis reformulações que favoreçam o desempenho de solvers de programação inteira. Tais melhorias podem resultar em modelos mais fortes e escaláveis, capazes de resolver instâncias maiores em tempos computacionais mais adequados.

% Também se espera que o desenvolvimento de abordagens heurísticas amplie as possibilidades de tratamento do PCmin em cenários nos quais métodos exatos se tornem inviáveis. A investigação de estratégias baseadas em metaheurísticas, como algoritmos genéticos e variantes em chaves aleatórias, poderá fornecer métodos capazes de produzir soluções de boa qualidade em tempo computacional reduzido, tornando o tratamento de instâncias maiores mais viável na prática.

% Adicionalmente, os resultados obtidos poderão servir de base para pesquisas futuras em problemas correlatos de precificação, como o PCmax e outras variantes estudadas na literatura. Assim, o projeto poderá contribuir tanto do ponto de vista teórico, ao aprofundar a compreensão estrutural do problema, quanto do ponto de vista computacional, ao propor e avaliar novas estratégias de resolução.

% \subsection{Desafios previstos}

% Entre os principais desafios previstos, destaca-se a própria complexidade computacional do PCmin, uma vez que se trata de um problema NP-difícil. Essa característica pode limitar a eficiência dos métodos exatos em instâncias de maior porte, mesmo quando são empregadas formulações fortalecidas. Dessa forma, é possível que os ganhos obtidos com melhorias de modelagem sejam graduais e dependam de sucessivos refinamentos.

% Outro desafio relevante está associado ao desenvolvimento e à calibração das heurísticas. A escolha de representações adequadas das soluções, bem como a definição de operadores de busca e parâmetros para metaheurísticas, pode exigir um processo iterativo de experimentação. Além disso, a comparação entre abordagens distintas demandará uma base experimental consistente, com instâncias representativas e critérios de avaliação bem definidos.

% Por fim, um desafio adicional consiste em equilibrar, dentro do tempo disponível para execução do projeto, a investigação de melhorias em modelos exatos e o desenvolvimento de heurísticas. Como ambas as frentes apresentam potencial de contribuição, será necessário conduzir o trabalho de forma criteriosa, priorizando as abordagens que se mostrarem mais promissoras ao longo da pesquisa.

Espera-se que o projeto contribua para o avanço do estudo do PCmin, ainda pouco explorado sob a perspectiva de formulações exatas e métodos de solução para instâncias de maior porte, como, por exemplo, aquelas com 1000 ou mais itens e compradores. 

Entre as contribuições esperadas, destacam-se o aperfeiçoamento das formulações matemáticas existentes, bem como o desenvolvimento de abordagens heurísticas capazes de produzir soluções de boa qualidade em tempo computacional reduzido quando métodos exatos se tornarem impraticáveis. Além disso, os resultados obtidos poderão servir de base para pesquisas futuras em variantes correlatas de problemas de precificação.

Entre os principais desafios previstos, sobressai a elevada complexidade computacional do PCmin, que pode limitar a efetividade dos métodos exatos mesmo com formulações fortalecidas. Também são esperadas dificuldades no desenvolvimento e na calibração das heurísticas, especialmente quanto à definição de representações, operadores de busca e parâmetros adequados. Por fim, será necessário equilibrar, dentro do tempo disponível, os esforços entre o aperfeiçoamento de modelos exatos e a investigação de abordagens heurísticas, priorizando as estratégias que se mostrarem mais promissoras ao longo da pesquisa.